\section{7.1 Calibration Quality Summary}
Intrinsics quality is strong across all cameras under current A4 ChArUco workflow, with sub-pixel to low-pixel reprojection errors. Extrinsics quality is non-uniform but usable after robust tag-map tuning and partial recalibration.

Representative final extrinsics RMSE values:
\begin{itemize}
  \item camNorth: 1.44 px
  \item camEast: 1.18 px
  \item camSouth: 5.23 px
  \item camWest: 2.26 px
\end{itemize}

The camSouth value reflects recent movement and remaining sensitivity.

\section{7.2 Ball Static GT: Raw Performance}
From \texttt{reports\_static\_raw/summary\_metrics.json}:
\begin{itemize}
  \item mean: 150.77 mm
  \item median: 156.55 mm
  \item RMSE: 167.39 mm
  \item P90: 236.38 mm
  \item P95: 288.34 mm
  \item max: 361.83 mm
  \item axis bias: ex +54.50 mm, ey +14.28 mm, ez -103.36 mm
  \item mean cameras used: 2.87
  \item mean reprojection: 6.01 px
\end{itemize}

Interpretation:
\begin{itemize}
  \item raw estimates are precise within windows (std low) but biased globally,
  \item Z-axis underestimation is dominant,
  \item camera coverage below 3 in many points limits triangulation robustness.
\end{itemize}

\section{7.3 Ball Static GT: Axis-Wise Correction}
After linear axis correction:
\begin{itemize}
  \item mean: 95.17 mm
  \item median: 84.18 mm
  \item RMSE: 102.23 mm
  \item P95: 166.51 mm
  \item max: 214.60 mm
\end{itemize}

Bias means become near zero by construction. Precision remains similar, showing correction mainly addresses systematic offset, not random jitter.

\subsection{Table 7.1 Ball static metrics (raw vs corrected)}
\begin{verbatim}
| Metric | Raw | Corrected |
|---|---:|---:|
| Mean (mm) | 150.77 | 95.17 |
| Median (mm) | 156.55 | 84.18 |
| RMSE (mm) | 167.39 | 102.23 |
| P90 (mm) | 236.38 | 142.18 |
| P95 (mm) | 288.34 | 166.51 |
| Max (mm) | 361.83 | 214.60 |
\end{verbatim}

\section{7.4 Dynamic Ball Behavior}
From dynamic summary:

\subsection{\texttt{ball\_slow}}
\begin{itemize}
  \item detect ratio: 0.975
  \item mean reproj: 4.03 px
  \item jump P95: 58.16 mm
  \item max jump: 173.07 mm
\end{itemize}

\subsection{\texttt{ball\_fast}}
\begin{itemize}
  \item detect ratio: 0.891
  \item mean reproj: 6.51 px
  \item jump P95: 462.70 mm
  \item max jump: 814.46 mm
\end{itemize}

\subsection{\texttt{no\_ball}}
\begin{itemize}
  \item detect ratio: 0.0 (desired)
\end{itemize}

\subsection{Table 7.2 Dynamic metrics summary}
\begin{verbatim}
| Clip | Detection ratio | Mean reproj (px) | P95 jump (mm) | Max jump (mm) |
|---|---:|---:|---:|---:|
| ball_slow | 0.975 | 4.03 | 58.16 | 173.07 |
| ball_fast | 0.891 | 6.51 | 462.70 | 814.46 |
| no_ball | 0.000 | - | - | - |
\end{verbatim}

Interpretation:
\begin{itemize}
  \item slow dynamics are acceptable,
  \item fast throws still produce outliers,
  \item false positives are strongly controlled.
\end{itemize}

\section{7.5 Joint Localization Results}
From \texttt{joint\_tuning\_20260310\_124311/reports/summary\_metrics.json}:
\begin{itemize}
  \item valid trials: 62 / 81
  \item mean error: 143.38 mm
  \item median: 148.90 mm
  \item RMSE: 147.73 mm
  \item P95: 198.73 mm
  \item max: 217.34 mm
  \item axis bias: ex +50.68 mm, ey +46.57 mm, ez -106.98 mm
\end{itemize}

Per-joint means:
\begin{itemize}
  \item left\_shoulder: 164.38 mm
  \item right\_hip: 150.38 mm
  \item right\_knee: 110.03 mm
\end{itemize}

\subsection{Table 7.3 Joint-touch metrics}
\begin{verbatim}
| Joint | Mean (mm) | RMSE (mm) | P95 (mm) |
|---|---:|---:|---:|
| left_shoulder | 164.38 | 165.60 | 199.54 |
| right_hip | 150.38 | 151.54 | 172.31 |
| right_knee | 110.03 | 118.68 | 170.75 |
\end{verbatim}

Interpretation:
\begin{itemize}
  \item shoulder is most difficult due to visibility and anatomical movement,
  \item knee is comparatively better,
  \item human placement uncertainty likely inflates measured error.
\end{itemize}

\section{7.6 Error Budget Decomposition}
Estimated contributors (qualitative ranking):
\begin{enumerate}
  \item camera movement and extrinsics drift (high),
  \item overlap and view geometry limitations (high),
  \item fast-motion blur (medium-high),
  \item person identity ambiguity (medium),
  \item intrinsics residual error (low-medium),
  \item manual GT placement uncertainty (medium, especially for joints).
\end{enumerate}

This ranking justifies prioritizing geometry discipline and overlap redesign before advanced model tuning.

\section{7.7 Why Skeleton Appears to "Lag" or "Jump"}
Observed lag/jumps were mainly caused by:
\begin{itemize}
  \item identity switching when a second person was visible,
  \item insufficient camera support for some joints in certain frames,
  \item sudden reprojection-consistent but semantically wrong triangulations under partial occlusion.
\end{itemize}

Mitigations used:
\begin{itemize}
  \item target person selection with area/confidence and temporal proximity,
  \item \texttt{pose-min-cams} gating,
  \item scene protocol with one active subject.
\end{itemize}

\section{7.8 Discussion: Is Correction Model Needed?}
Correction model is useful only when:
\begin{itemize}
  \item calibration is reasonably stable,
  \item residual error is predominantly systematic,
  \item and corrected model is applied within calibrated operating volume.
\end{itemize}

If camera positions change, correction learned from previous geometry may become invalid. Thus correction is a secondary layer after geometry stabilization.

\section{7.9 Readiness Conclusion from Results}
The perception stack is suitable for:
\begin{itemize}
  \item visualization and analytics,
  \item target-zone event prototyping,
  \item supervised launcher integration experiments.
\end{itemize}

It is not yet sufficient for unsupervised autonomous actuation against strict centimeter-level body-part targets.


\section{7.10 Interpretation for Stakeholders}
From a stakeholder perspective, the key message is:
\begin{itemize}
  \item the system is already useful for visual analytics and controlled targeting research,
  \item the system is not yet validated for unsupervised high-precision body-part launch commands.
\end{itemize}

This balanced interpretation protects project credibility and supports realistic planning.

\section{7.11 What Improvements Give Highest Return}
Based on measured data, highest-return improvements are:
\begin{enumerate}
  \item increase average camera participation from ~2.87 toward >=3.3,
  \item harden camSouth (or weakest camera) extrinsics and overlap,
  \item maintain strict single-subject protocol during pose tracking,
  \item use rigid calibration targets for final controller tuning.
\end{enumerate}

Detector model changes are secondary until these geometric constraints are addressed.

\section{7.12 Recommended Acceptance Gates Before Actuation}
Suggested acceptance gates before any autonomous firing trial:
\begin{itemize}
  \item static corrected mean <= 80 mm,
  \item static corrected P95 <= 130 mm,
  \item dynamic fast outlier max <= 500 mm,
  \item no-ball false positive ratio <= 0.5\%,
  \item joint-touch median <= 120 mm in central operational zone.
\end{itemize}

These thresholds can be refined, but explicit gates are necessary for safe progression.

\section{7.13 Evidence Strength and Remaining Uncertainty}
Evidence strength is high for:
\begin{itemize}
  \item pipeline reproducibility,
  \item static bias characterization,
  \item dynamic failure mode identification.
\end{itemize}

Remaining uncertainty is high for:
\begin{itemize}
  \item full-volume joint accuracy under free movement,
  \item controller-level hit accuracy once launcher dynamics are integrated.
\end{itemize}

This uncertainty profile is acceptable for concluding perception-stage success and moving to HIL-stage research.



\section{7.14 Comparative Interpretation: Ball vs Joint Error}
Ball static estimates benefit from a clear visual target and rigid center concept. Joint estimates depend on semantic body landmarks and pose model interpretation. Therefore, it is expected that joint errors are larger and less uniform.

This comparison is useful when setting performance expectations for control logic: joint-targeted commands should carry stricter confidence and safety gating than zone-level commands.

\section{7.15 Spatial Distribution of Errors}
Although aggregate metrics are informative, spatial distribution matters. Errors were generally higher near weaker-overlap regions and for heights with reduced multi-camera support. This suggests future target-selection logic should consider "safe operational zones" where geometry is best conditioned.

\section{7.16 Precision vs Accuracy Distinction in This Project}
The project repeatedly observed low within-window jitter (precision) with nontrivial absolute bias (accuracy issue). This distinction is essential:
\begin{itemize}
  \item precision indicates stable estimations,
  \item bias indicates frame alignment/calibration issues.
\end{itemize}

Correction models improved accuracy but do not change underlying precision characteristics.

\section{7.17 Implications for Commanded Target Types}
Given current error profile, target categories can be staged:
\begin{itemize}
  \item ready now: larger zones (1x1 or 1x1.5 m planes), body-center approximations,
  \item near-term: lower-body targets with better current metrics,
  \item later: small body-part targets (e.g., shoulder points) requiring tighter error bounds.
\end{itemize}

This staged approach reduces risk while preserving project momentum.

\section{7.18 Statistical Confidence and Reporting Transparency}
All reported metrics are deterministic from stored trial data, and missing/failed trials are explicitly reported rather than excluded silently. This transparency is important for committee trust and for future replication.



\section{7.19 Scenario-Based Interpretation for Training Use Cases}
To connect metrics with training decisions, consider three scenarios:

\begin{enumerate}
  \item Large-zone passing drills (>=1 m targets): current corrected ball errors are often acceptable, especially in central volume.
  \item Body-region targeting (torso/hip zones): feasible with strict confidence gating and supervised operation.
  \item Small body-part targeting (e.g., shoulder point): currently high risk without further geometry improvement.
\end{enumerate}

This scenario framing helps stakeholders make informed deployment decisions.

\section{7.20 Case Study: Camera Movement and Recovery}
When camSouth moved, pipeline accuracy degraded. Partial recalibration and merge restored operational consistency without redoing full system setup. This case demonstrates practical maintainability of the chosen architecture and validates the partial-recalibration strategy.

\section{7.21 Case Study: Two-Person Scene Instability}
When a second person remained visible, skeleton identity instability increased. Enforcing single-subject protocol resolved major artifacts. This case emphasizes that protocol constraints are not optional for meaningful evaluation.

\section{7.22 Case Study: Ball Visibility and Camera Count}
Trials with only one or two camera views were significantly more unstable. Redesigning trial points toward higher overlap improved reliability. This reinforces observability-first design for future protocols and control logic.

\section{7.23 Practical Performance Envelope}
Given current results, practical envelope is:
\begin{itemize}
  \item best performance in central region with >=3 camera support,
  \item reduced reliability near edge/corner zones,
  \item higher instability in high-speed throws and partial occlusions.
\end{itemize}

Controllers should incorporate this envelope and restrict autonomous behavior outside it.

\section{7.24 Chapter Summary}
Results confirm meaningful progress and clear remaining work. The key conclusion is that geometric and protocol discipline produce stronger gains than isolated detector tuning at this stage.



\section{7.25 Multi-Metric Decision Policy for Next-Step Progression}
A single metric cannot determine readiness. A practical policy should combine:
\begin{itemize}
  \item corrected static mean and P95,
  \item dynamic outlier rate,
  \item no-ball false-positive rate,
  \item camera participation statistics,
  \item calibration health state.
\end{itemize}

Only if all indicators are within thresholds should stage transition be approved.

\section{7.26 Why P95 Is More Informative Than Mean for Safety}
Mean error can hide dangerous tails. P95 better reflects near-worst operational behavior without being dominated by one extreme outlier. For launcher safety, tail behavior is critical because one extreme miss can have severe consequences.

\section{7.27 Handling Discrepancy Between Visual Plausibility and Metrics}
There were cases where rendered trajectories looked acceptable but metric errors remained high. This discrepancy arises because visualization perspective can mask absolute displacement. Therefore, quantitative reporting must take precedence in technical decisions.

\section{7.28 Error Reduction Prioritization Plan}
Ordered plan for further error reduction:
\begin{enumerate}
  \item lock hardware mounts and tag planarity,
  \item maximize overlap and tag observability,
  \item rerun full rigid-point GT,
  \item recalibrate correction model only after steps 1-3,
  \item tune dynamic filters for fast motion.
\end{enumerate}

This sequence avoids premature algorithm tuning on unstable geometry.

\section{7.29 Lessons for Thesis Defense}
The strongest defense narrative is:
\begin{itemize}
  \item clear problem framing,
  \item measurable protocol execution,
  \item transparent limitations,
  \item concrete next-step gates.
\end{itemize}

The current results satisfy this narrative and show rigorous research process.
